%% ma: L12-B14-0131
%% lop: 12
%% chuong: 5
%% bai: 14
%% dang: 12.14.39
%% loai: TN4
%% muc: VD
%% dap_an: "A"
%% nguon: "Ảnh giáo viên gửi 25/11/2026 (tài liệu ôn tập Phương trình mặt phẳng), Phần I Câu 22"
%% kien_thuc: "Bài 14 (phương trình mặt phẳng); hình chóp có SA vuông góc với đáy, đáy là hình vuông (BD⊥(SAC))"
%% trang_thai: cho-duyet
%% ly_do: "Dạng toán mới đề xuất, chờ giáo viên duyệt"
\begin{ex}
Trong không gian $Oxyz$, cho hình chóp $S.ABCD$ có đáy là hình vuông và $SA$ vuông góc với đáy. Cho biết $B(2;3;7)$, $D(4;1;3)$. Lập phương trình mặt phẳng $(SAC)$.
\choice{\True $x-y-2z+9=0$}{$x-y+2z+9=0$}{$x-y-2z-9=0$}{$x+y-2z+9=0$}
\loigiai{Vì $ABCD$ là hình vuông nên $BD\perp AC$; lại có $BD\perp SA$, nên $BD\perp(SAC)$. Vậy $\overrightarrow{BD}=(2;-2;-4)=2(1;-1;-2)$ là một vectơ pháp tuyến của $(SAC)$.
Mặt phẳng $(SAC)$ đi qua tâm hình vuông là trung điểm $I(3;2;5)$ của $BD$: $(x-3)-(y-2)-2(z-5)=0\Leftrightarrow x-y-2z+9=0$.}
\end{ex}
